%LTeX: language=fr-FR
\documentclass[12pt,a4paper]{article}

% ---------- Langue et encodage ----------
\usepackage[french]{babel}
%\usepackage[T1]{fontenc}
%\usepackage[utf8]{inputenc}

% ---------- Mise en page ----------
\usepackage[margin=2.2cm]{geometry}
\usepackage{array,booktabs,tabularx}
\usepackage{amsmath,amssymb}
\usepackage{xcolor}
\usepackage{tcolorbox}
\tcbuselibrary{skins,breakable}
\usepackage{enumitem}
\usepackage{graphicx}
\usepackage{tikz}
\usepackage{pgfplots}
\pgfplotsset{compat=1.18}
\usepackage{fancyhdr}
\usepackage{makecell}
\usepackage{pgffor}
\usepackage{currfile}
\usepackage{lastpage}

\graphicspath{ {../../Images/} }
\input{\currfiledir ../def.tex}
\def \Document {Activité 2}


% ---------- Couleurs ----------
\definecolor{bleuF}{RGB}{47,87,163}   
\definecolor{orH}{RGB}{242,170,60}    
\definecolor{cadre}{RGB}{47,87,163}

% ---------- Boîtes ----------
\newtcolorbox{objectif}{
  colback=bleuF!6, colframe=cadre, boxrule=0.6pt, arc=2pt,
  left=6pt, right=6pt, top=4pt, bottom=4pt,
  title={\textbf{Objectif}}, fonttitle=\bfseries,
  coltitle=cadre, colbacktitle=bleuF!12
}
\newtcolorbox{bilan}{
  colback=orH!8, colframe=cadre, boxrule=0.6pt, arc=2pt,
  left=6pt, right=6pt, top=4pt, bottom=4pt,
  title={\textbf{Bilan}}, fonttitle=\bfseries,
  coltitle=cadre, colbacktitle=orH!20,
  breakable
}
\newtcolorbox{correctionbox}{
  colback=green!5, colframe=green!40!black, boxrule=0.6pt, arc=2pt,
  left=6pt, right=6pt, top=4pt, bottom=4pt,
  breakable
}

% \checkmark symbol fallback (amssymb provides \checkmark)

% ---------- Ligne de réponse pointillée ----------
\newcommand{\reponse}[1]{%
  \par\vspace{2mm}
  \foreach \i in {1,...,#1}{\rule{0pt}{0pt}\dotfill\\[3mm]}%
}

\begin{document}


\pagestyle{fancy}

\fancyhf{} % clear existing header/footer entries
\fancyhead[R]{ {\small \Niveau}}
\fancyhead[L]{ {\small \ChapitreCourt}}
\fancyhead[C]{ {\small \Document{} -- Corrigé}}
%\fancyhead[R]{\large \textbf{\textbf{\Document\ (B) -- 1h}}}
\fancyfoot[L]{ A. Fernandez }
\fancyfoot[R]{{\small \thepage\  / \pageref{LastPage}}}

\setlength{\parindent}{0pt}
\newcommand{\rep}[1]{\textcolor{cadre}{\textbf{#1}}}

\section*{Le budget des repas}

\begin{enumerate}[label=\textbf{\arabic*.}, leftmargin=1.3em]
\item
\begin{correctionbox}
Chaque année, on ajoute $6$ enfants : $u(n+1) = u(n) + \rep{6}$.
\end{correctionbox}

\item
\begin{correctionbox}
$10\,\%$ de $100$~€, c'est $10$~€. Le nouveau prix est $100 + 10 = 110$~€, et $110 = 1{,}1 \times 100$.

Augmenter une quantité de $10\,\%$, c'est la multiplier par \rep{$1{,}1$} $\left(= 1 + \dfrac{10}{100}\right)$.
\end{correctionbox}

\item
\begin{correctionbox}
$v(1) = 1{,}1 \times v(0) = 1{,}1 \times 40\,000 = 44\,000$.

$v(2) = 1{,}1 \times v(1) = 1{,}1 \times 44\,000 = 48\,400$.

$v(3) = 1{,}1 \times v(2) = 1{,}1 \times 48\,400 = 53\,240$.

$v(4) = 1{,}1 \times v(3) = 1{,}1 \times 53\,240 = 58\,564$.
\end{correctionbox}
\begin{center}
\renewcommand{\arraystretch}{1.8}
\begin{tabular}{|c|*{5}{>{\centering\arraybackslash}p{1.6cm}|}}
\hline
$n$ & $0$ & $1$ & $2$ & $3$ & $4$ \\
\hline
$v(n)$ & $40\,000$ & \rep{44\,000} & \rep{48\,400} & \rep{53\,240} & \rep{58\,564} \\
\hline
\end{tabular}
\end{center}
Chaque année, le budget est multiplié par $1{,}1$ : $v(n+1) = \rep{1{,}1} \times v(n)$.

\item
\begin{correctionbox}
Non. Pour calculer $v(10)$, il faut connaître $v(9)$ ; pour calculer $v(9)$, il faut connaître $v(8)$, etc. Il faut donc calculer tous les termes un par un, de $v(1)$ jusqu'à $v(10)$.

Pour $u(10)$, la formule $u(n) = 20 + 6n$ donne le résultat directement.
\end{correctionbox}
\end{enumerate}

\begin{bilan}
\begin{itemize}[leftmargin=1.2em]
\item \textbf{Relation de récurrence} : on calcule un terme à partir du terme \rep{précédent}. Il faut connaître le premier terme. Exemples : $u(0) = 20$ et $u(n+1) = \rep{$u(n) + 6$}$ ; $v(0) = \rep{$40\,000$}$ et $v(n+1) = \rep{$1{,}1 \times v(n)$}$.
\item Avec une formule explicite, on calcule directement n'importe quel terme. Avec une relation de récurrence, il faut calculer \rep{tous} les termes précédents.
\end{itemize}
\end{bilan}

\end{document}
